% =============================================================================
% 3. SOLUCIÓN
% =============================================================================

\section{Solución}

\subsection{Sistema}

\subsubsection{Descripción general}

Smart-Park es una plataforma web de arquitectura tipo marketplace, con diseño completamente responsivo para computadoras, tablets y teléfonos móviles, instalable como aplicación (PWA). La arquitectura se organiza en tres niveles jerárquicos que separan la administración central, la operación de cada local afiliado y el acceso de los distintos perfiles de usuario:

\begin{itemize}
    \item \textbf{Nivel 1 --- Plataforma Smart-Park (nivel central).} Está desplegada en la nube y concentra la lógica de negocio común a toda la red:
    \begin{itemize}
        \item Autenticación, control de roles y seguridad perimetral.
        \item Gestión de usuarios, afiliaciones y estacionamientos.
        \item Reservas, tarifas y cobros de toda la red.
        \item Pagos electrónicos y liquidaciones financieras.
        \item Motor de visión artificial y escaneo periódico de las cámaras.
        \item Procesos automáticos (liberación de reservas vencidas y copias de seguridad).
        \item Difusión de eventos en tiempo real y mensajería asíncrona.
        \item Reportes consolidados, auditoría y configuración global.
        \item Evaluación técnica de los locales para habilitar la detección de vehículos con inteligencia artificial.
    \end{itemize}
    \item \textbf{Nivel 2 --- Nivel de estacionamiento (operación de cada local).} Corresponde a la infraestructura y la operación física de cada local afiliado:
    \begin{itemize}
        \item El plano digital del local, con sus espacios y su estado.
        \item Las cámaras IP instaladas en el local (en los locales habilitados), a las que la plataforma central accede para obtener las imágenes.
        \item La estación de garita, que es el navegador del trabajador, desde donde se validan los pases QR y se registran los ingresos y salidas.
        \item La gestión local de tarifas, personal, cámaras e incidencias, a cargo del administrador del local.
    \end{itemize}
    \item \textbf{Nivel 3 --- Usuarios.} Comprende los cuatro perfiles definidos (administrador de plataforma, administrador de local, trabajador del local y conductor). Todos acceden mediante interfaces web adaptadas a sus permisos, bajo un esquema de Control de Acceso Basado en Roles (RBAC).
\end{itemize}

\subsubsection{Características principales}

Las características que distinguen a Smart-Park como plataforma de estacionamiento inteligente son:

\begin{table}[H]
\caption{Características Principales del Sistema}
\begin{tabularx}{\textwidth}{p{3.5cm} X}
\toprule
\textbf{Característica} & \textbf{Descripción} \\
\midrule
Marketplace multi-tenant & Conecta múltiples estacionamientos afiliados con conductores, cada local opera con autonomía sobre sus datos, tarifas y políticas. \\
\midrule
Disponibilidad en tiempo real & Actualización inmediata del estado de espacios mediante WebSocket (WSS), con difusión selectiva por local y por usuario. \\
\midrule
Editor de planos 2D (CAD) & Permite al administrador de local dibujar la distribución física del estacionamiento: espacios, muros, accesos, garita y pasos peatonales. \\
\midrule
Reservas flexibles & Tres modalidades: estándar (horas/minutos), anticipada (fecha futura) y abono (semanal, mensual o fraccionado). \\
\midrule
Pase QR digital & Código QR con temporizador de tolerancia, verificable en una página pública de validación por el personal de garita. \\
\midrule
Visión artificial & Detección de vehículos con YOLOv8 y OpenCV, habilitada tras evaluación técnica de cámaras, iluminación y red del local. \\
\midrule
Pagos electrónicos & Integración con Culqi (tarjetas, Yape, Plin) y PayPal, con cálculo de comisiones y liquidación automática. \\
\midrule
Seguridad integral & RBAC, JWT en cookies HttpOnly, autenticación con Google, PIN administrativo, limitación de intentos y bitácora de auditoría. \\
\midrule
PWA responsiva & Interfaz web instalable en cualquier dispositivo, adaptada a los cuatro perfiles de usuario. \\
\midrule
Geolocalización y mapas & Búsqueda de estacionamientos mediante Leaflet con capas de calles y satélite, filtros por zona, tipo y tarifa. \\
\bottomrule
\end{tabularx}
\end{table}

\subsubsection{Componentes del sistema (contenedores)}

La plataforma se compone de los siguientes contenedores lógicos, cada uno con una responsabilidad específica dentro de la arquitectura:

\begin{table}[H]
\caption{Componentes (Contenedores) del Sistema}
\begin{tabularx}{\textwidth}{p{3.5cm} p{2.5cm} X}
\toprule
\textbf{Contenedor} & \textbf{Tecnología} & \textbf{Responsabilidad} \\
\midrule
Aplicación Web (SPA) & React 19 + Vite & Interfaz de usuario responsiva, enrutamiento del lado del cliente, comunicación con la API mediante HTTPS y con el servidor de WebSocket mediante WSS. Instalable como PWA. \\
\midrule
API central (Backend) & FastAPI + SQLAlchemy asíncrono & Lógica de negocio, autenticación, autorización, gestión de reservas, pagos, afiliaciones, visión artificial, reportes y auditoría. Expone endpoints REST con JSON. \\
\midrule
Servidor de WebSocket & FastAPI (módulo integrado) & Difusión en tiempo real del estado de los espacios, notificaciones instantáneas y eventos de reserva. \\
\midrule
Motor de visión artificial & Python + YOLOv8 + OpenCV & Captura periódica de imágenes desde las cámaras IP, detección de vehículos y clasificación del estado de los espacios. \\
\midrule
Cola de mensajes & RabbitMQ (CloudAMQP) & Desacopla la lógica de negocio de los procesos asíncronos: notificaciones, liberación de reservas, copias de seguridad y eventos de cámara. \\
\midrule
Caché en memoria & Redis (Railway) & Almacena sesiones, tokens de acceso, estado temporal de espacios y resultados de consultas frecuentes. \\
\midrule
Base de datos & PostgreSQL (Railway) & Almacenamiento persistente de usuarios, estacionamientos, espacios, reservas, pagos, reseñas, incidencias, auditoría y configuración. \\
\midrule
Almacenamiento de archivos & Volumen \texttt{/data} (Railway) & Fotografías de vehículos, evidencias de incidencias, logotipos de locales y copias de seguridad de la base de datos. \\
\bottomrule
\end{tabularx}
\end{table}

\subsection{Diagrama de Infraestructura}

El diagrama de infraestructura (Figura~2) muestra la distribución física de los componentes del sistema, sus tecnologías base y los protocolos de comunicación entre ellos.

% Aquí debes incluir tu imagen del diagrama de infraestructura:
% \begin{figure}[H]
% \caption{Diagrama de Infraestructura del Sistema Smart-Park}
% \centering
% \includegraphics[width=0.95\textwidth]{imagenes/diagrama_infraestructura.png}
% \end{figure}
% \textit{Nota.} Elaboración propia.

La infraestructura de producción se despliega en Railway con los siguientes servicios:

\begin{itemize}
    \item \textbf{Servicio principal:} contenedor Docker multi-stage que ejecuta la API (FastAPI con Uvicorn), el servidor WebSocket y el motor de visión artificial. Se expone mediante el dominio personalizado con certificado TLS.
    \item \textbf{PostgreSQL:} instancia gestionada por Railway con respaldo automático diario y cifrado en reposo.
    \item \textbf{Redis:} instancia gestionada por Railway, utilizada como caché de sesiones y estado temporal de espacios.
    \item \textbf{RabbitMQ:} instancia en CloudAMQP (plan gestionado), conectada al servicio principal mediante AMQPS (AMQP cifrado).
    \item \textbf{Volumen persistente:} montado en \texttt{/data} para almacenar fotografías y copias de seguridad.
    \item \textbf{CDN y frontend:} la aplicación React se sirve mediante Vercel, conectándose a la API central por HTTPS.
\end{itemize}
